
%%%%%%%%%%%%%%%%%%%%%%%%%%%%%%%%%%%%%%%%%%%%%%%%%%%%%%%%%%%%%%%%%%%%%%%%%%%%%%%%
\chapter{The Origin of Stochasticity in Physical Theories}
\label{ch:OriginStochasticity}
%%%%%%%%%%%%%%%%%%%%%%%%%%%%%%%%%%%%%%%%%%%%%%%%%%%%%%%%%%%%%%%%%%%%%%%%%%%%%%%%

%%%%%%%%%%%%%%%%%%%%%%%%%%%%%%%%%%%%%%%%%%%%%%%%%%%%%%%%%%%%%%%%%%%%%%%%%%%%%%%%
\section{Introduction}
%%%%%%%%%%%%%%%%%%%%%%%%%%%%%%%%%%%%%%%%%%%%%%%%%%%%%%%%%%%%%%%%%%%%%%%%%%%%%%%%

Stochastic equations appear throughout modern physics.

Examples include Brownian motion, kinetic theory, stochastic
electrodynamics, open quantum systems, fluctuating hydrodynamics,
stochastic mechanics, and quantum measurement theory.

Although these theories employ mathematically similar stochastic
equations, the physical origin of the stochasticity differs
fundamentally.

Failure to distinguish these origins has often led to conceptual
confusion.

The purpose of this chapter is to classify the principal mechanisms by
which stochastic behavior emerges in physical theories and to place
stochastic electrodynamics within this broader framework.

%%%%%%%%%%%%%%%%%%%%%%%%%%%%%%%%%%%%%%%%%%%%%%%%%%%%%%%%%%%%%%%%%%%%%%%%%%%%%%%%
\section{Fundamental Classification}
%%%%%%%%%%%%%%%%%%%%%%%%%%%%%%%%%%%%%%%%%%%%%%%%%%%%%%%%%%%%%%%%%%%%%%%%%%%%%%%%

At the highest level, stochasticity may arise in four conceptually
distinct ways.

\begin{enumerate}

\item
No stochasticity: the microscopic equations are deterministic and all
degrees of freedom are retained.

\item
Effective stochasticity: deterministic microscopic dynamics is reduced
by eliminating unobserved variables.

\item
Phenomenological stochasticity: random forces are introduced as an
effective description without specifying microscopic origins.

\item
Fundamental stochasticity: randomness is postulated as an irreducible
feature of the theory.

\end{enumerate}

This classification provides a common language for comparing apparently
very different physical theories.

%%%%%%%%%%%%%%%%%%%%%%%%%%%%%%%%%%%%%%%%%%%%%%%%%%%%%%%%%%%%%%%%%%%%%%%%%%%%%%%%
\section{Deterministic Hamiltonian Dynamics}
%%%%%%%%%%%%%%%%%%%%%%%%%%%%%%%%%%%%%%%%%%%%%%%%%%%%%%%%%%%%%%%%%%%%%%%%%%%%%%%%

Classical mechanics represents the simplest case.

The evolution is generated by Hamilton's equations,

\begin{align}
\dot q^i
&=
\frac{\partial H}{\partial p_i},
\\
\dot p_i
&=
-
\frac{\partial H}{\partial q^i}.
\end{align}

No randomness appears.

Any statistical description reflects only incomplete knowledge of
initial conditions rather than stochastic dynamics.

%%%%%%%%%%%%%%%%%%%%%%%%%%%%%%%%%%%%%%%%%%%%%%%%%%%%%%%%%%%%%%%%%%%%%%%%%%%%%%%%
\section{Reduction of Hidden Degrees of Freedom}
%%%%%%%%%%%%%%%%%%%%%%%%%%%%%%%%%%%%%%%%%%%%%%%%%%%%%%%%%%%%%%%%%%%%%%%%%%%%%%%%

Many stochastic equations originate from deterministic systems by
eliminating unobserved variables.

Suppose the microscopic state is

\[
X=(x,y),
\]

where

\[
x
\]

denotes observed variables and

\[
y
\]

represents inaccessible environmental degrees of freedom.

Projection onto

\[
x
\]

produces equations of the form

\begin{equation}
\dot x
=
F(x)
+
\int_0^t
K(t-s)x(s)\,ds
+
\xi(t),
\end{equation}

where

\[
K(t)
\]

is a memory kernel and

\[
\xi(t)
\]

is an effective fluctuating force.

Generalized Langevin equations, Mori--Zwanzig theory, and stochastic
electrodynamics all belong to this class.

%%%%%%%%%%%%%%%%%%%%%%%%%%%%%%%%%%%%%%%%%%%%%%%%%%%%%%%%%%%%%%%%%%%%%%%%%%%%%%%%
\section{Phenomenological Noise}
%%%%%%%%%%%%%%%%%%%%%%%%%%%%%%%%%%%%%%%%%%%%%%%%%%%%%%%%%%%%%%%%%%%%%%%%%%%%%%%%

In many applications, the microscopic origin of the random force is
irrelevant.

One simply writes

\begin{equation}
m\ddot x
=
F(x)
+
\xi(t),
\end{equation}

where

\[
\xi(t)
\]

is prescribed statistically.

This description underlies much of engineering, chemical physics,
population dynamics and turbulence modeling.

Its predictive success does not depend upon identifying the microscopic
source of the fluctuations.

%%%%%%%%%%%%%%%%%%%%%%%%%%%%%%%%%%%%%%%%%%%%%%%%%%%%%%%%%%%%%%%%%%%%%%%%%%%%%%%%
\section{Fundamental Stochastic Theories}
%%%%%%%%%%%%%%%%%%%%%%%%%%%%%%%%%%%%%%%%%%%%%%%%%%%%%%%%%%%%%%%%%%%%%%%%%%%%%%%%

Some approaches regard stochasticity as an intrinsic property of
nature.

Examples include

\begin{itemize}

\item
Nelson's stochastic mechanics,

\item
Continuous spontaneous localization,

\item
Certain dynamical collapse models.

\end{itemize}

In these theories,

the stochastic process is not derived from hidden variables but is
postulated as a basic element of the dynamics.

%%%%%%%%%%%%%%%%%%%%%%%%%%%%%%%%%%%%%%%%%%%%%%%%%%%%%%%%%%%%%%%%%%%%%%%%%%%%%%%%
\section{Position of Stochastic Electrodynamics}
%%%%%%%%%%%%%%%%%%%%%%%%%%%%%%%%%%%%%%%%%%%%%%%%%%%%%%%%%%%%%%%%%%%%%%%%%%%%%%%%

Within this classification,

stochastic electrodynamics occupies an intermediate position.

The microscopic particle--field dynamics is deterministic.

Randomness enters only after the electromagnetic field is treated
statistically.

Consequently,

SED is neither an intrinsically stochastic theory nor merely a
phenomenological Langevin model.

Instead,

it is an example of reduced deterministic dynamics.

%%%%%%%%%%%%%%%%%%%%%%%%%%%%%%%%%%%%%%%%%%%%%%%%%%%%%%%%%%%%%%%%%%%%%%%%%%%%%%%%
\section{Comparison with Nelson Mechanics}
%%%%%%%%%%%%%%%%%%%%%%%%%%%%%%%%%%%%%%%%%%%%%%%%%%%%%%%%%%%%%%%%%%%%%%%%%%%%%%%%

The contrast with Nelson's theory is instructive.

In Nelson mechanics,

Brownian motion is assumed from the outset.

The Wiener process is fundamental.

In stochastic electrodynamics,

the fluctuating electromagnetic field has its own deterministic
equations of motion.

The random force acting on the particle emerges only after averaging
over unknown field configurations.

The mathematical equations may appear similar,

yet the conceptual foundations are entirely different.

%%%%%%%%%%%%%%%%%%%%%%%%%%%%%%%%%%%%%%%%%%%%%%%%%%%%%%%%%%%%%%%%%%%%%%%%%%%%%%%%
\section{Implications for SHP}
%%%%%%%%%%%%%%%%%%%%%%%%%%%%%%%%%%%%%%%%%%%%%%%%%%%%%%%%%%%%%%%%%%%%%%%%%%%%%%%%

The SHP framework raises an important question.

Should stochasticity be introduced as

\begin{enumerate}

\item
a fundamental postulate,

\item
an effective consequence of eliminating hidden degrees of freedom,

\item
or an emergent property of covariant many-body dynamics?

The answer determines not only the mathematical formulation of the
theory but also its physical interpretation.

Throughout the remainder of this monograph we shall argue that the third
possibility deserves particularly careful investigation.

%%%%%%%%%%%%%%%%%%%%%%%%%%%%%%%%%%%%%%%%%%%%%%%%%%%%%%%%%%%%%%%%%%%%%%%%%%%%%%%%
\section*{Critical Assessment}
%%%%%%%%%%%%%%%%%%%%%%%%%%%%%%%%%%%%%%%%%%%%%%%%%%%%%%%%%%%%%%%%%%%%%%%%%%%%%%%%

The classification presented in this chapter emphasizes that stochastic
equations alone do not determine the physical interpretation of a
theory. Similar Langevin or Fokker--Planck equations may arise from
fundamentally different mechanisms, including coarse-graining,
statistical uncertainty, interaction with an environment, or genuinely
stochastic microscopic laws.

Recognizing these distinctions is essential when comparing stochastic
electrodynamics, Nelson mechanics, generalized Langevin theory, and
future stochastic extensions of the SHP framework. The mathematical
similarities among these theories are substantial, but the conceptual
origin of their stochasticity is not.
%%%%%%%%%%%%%%%%%%%%%%%%%%%%%%%%%%%%%%%%%%%%%%%%%%%%%%%%%%%%%%%%%%%%%%%%%%%%%%%%
\chapter{Axioms of Covariant Stochastic Hamiltonian Dynamics}
\label{ch:CSHTAxioms}
%%%%%%%%%%%%%%%%%%%%%%%%%%%%%%%%%%%%%%%%%%%%%%%%%%%%%%%%%%%%%%%%%%%%%%%%%%%%%%%%

%%%%%%%%%%%%%%%%%%%%%%%%%%%%%%%%%%%%%%%%%%%%%%%%%%%%%%%%%%%%%%%%%%%%%%%%%%%%%%%%
\section{Introduction}
%%%%%%%%%%%%%%%%%%%%%%%%%%%%%%%%%%%%%%%%%%%%%%%%%%%%%%%%%%%%%%%%%%%%%%%%%%%%%%%%

The preceding chapters have demonstrated that stochasticity enters
physical theories in several distinct ways. Brownian motion,
generalized Langevin equations, stochastic electrodynamics, and
stochastic mechanics all employ stochastic differential equations, but
their physical interpretation differs fundamentally.

Our objective is now different.

Rather than reviewing existing theories, we seek a general framework
within which relativistic stochastic dynamics can be formulated without
abandoning the Hamiltonian structure that underlies both classical
mechanics and the Stueckelberg--Horwitz--Piron (SHP) theory.

The philosophy adopted here is conservative.

We assume that deterministic covariant Hamiltonian dynamics constitutes
the fundamental description. Stochasticity appears only after reducing
or averaging over unresolved degrees of freedom.

%%%%%%%%%%%%%%%%%%%%%%%%%%%%%%%%%%%%%%%%%%%%%%%%%%%%%%%%%%%%%%%%%%%%%%%%%%%%%%%%
\section{Requirements}
%%%%%%%%%%%%%%%%%%%%%%%%%%%%%%%%%%%%%%%%%%%%%%%%%%%%%%%%%%%%%%%%%%%%%%%%%%%%%%%%

Any acceptable relativistic stochastic dynamics should satisfy the
following fundamental requirements.

\subsection*{Requirement 1: Lorentz Covariance}

The equations of motion must transform covariantly under the Poincaré
group.

No preferred inertial frame should appear in the fundamental equations.

%%%%%%%%%%%%%%%%%%%%%%%%%%%%%%%%%%%%%%%%%%%%%%%%%%%%%%%%%%%%%%%%%%%%%%%%%%%%%%%%

\subsection*{Requirement 2: Covariant Evolution Parameter}

Evolution should be parametrized by an invariant parameter

\[
\tau,
\]

rather than by the coordinate time of a particular observer.

This requirement is naturally satisfied within the SHP framework.

%%%%%%%%%%%%%%%%%%%%%%%%%%%%%%%%%%%%%%%%%%%%%%%%%%%%%%%%%%%%%%%%%%%%%%%%%%%%%%%%

\subsection*{Requirement 3: Hamiltonian Structure}

In the absence of stochastic interactions the equations must reduce to

\begin{align}
\frac{dx^\mu}{d\tau}
&=
\frac{\partial K}{\partial p_\mu},
\\
\frac{dp_\mu}{d\tau}
&=
-
\frac{\partial K}{\partial x^\mu},
\end{align}

where

\[
K(x,p)
\]

is the SHP Hamiltonian.

%%%%%%%%%%%%%%%%%%%%%%%%%%%%%%%%%%%%%%%%%%%%%%%%%%%%%%%%%%%%%%%%%%%%%%%%%%%%%%%%

\subsection*{Requirement 4: Causality}

The stochastic evolution at parameter value

\[
\tau
\]

may depend only upon earlier values

\[
\tau'<\tau.
\]

Memory kernels must therefore satisfy

\[
K(\tau-\tau')
=
0,
\qquad
\tau'<\tau.
\]

%%%%%%%%%%%%%%%%%%%%%%%%%%%%%%%%%%%%%%%%%%%%%%%%%%%%%%%%%%%%%%%%%%%%%%%%%%%%%%%%

\subsection*{Requirement 5: Covariant Probability Conservation}

The probability density

\[
\rho(x,p,\tau)
\]

must remain normalized,

\[
\int
\rho
\,d\Gamma
=
1,
\]

where

\[
d\Gamma
\]

denotes the invariant phase-space measure.

%%%%%%%%%%%%%%%%%%%%%%%%%%%%%%%%%%%%%%%%%%%%%%%%%%%%%%%%%%%%%%%%%%%%%%%%%%%%%%%%

\subsection*{Requirement 6: Deterministic Limit}

In the absence of stochastic interactions,

\[
\xi^\mu
\rightarrow
0,
\]

the theory must reduce continuously to ordinary SHP dynamics.

%%%%%%%%%%%%%%%%%%%%%%%%%%%%%%%%%%%%%%%%%%%%%%%%%%%%%%%%%%%%%%%%%%%%%%%%%%%%%%%%

\subsection*{Requirement 7: Nonrelativistic Limit}

For

\[
v\ll c,
\]

the theory should reproduce the known equations of stochastic mechanics
or generalized Langevin dynamics, depending on the physical mechanism
used to generate the stochastic force.

%%%%%%%%%%%%%%%%%%%%%%%%%%%%%%%%%%%%%%%%%%%%%%%%%%%%%%%%%%%%%%%%%%%%%%%%%%%%%%%%
\section{General Form of the Equations}
%%%%%%%%%%%%%%%%%%%%%%%%%%%%%%%%%%%%%%%%%%%%%%%%%%%%%%%%%%%%%%%%%%%%%%%%%%%%%%%%

The preceding requirements suggest equations of the form

\begin{align}
\frac{dx^\mu}{d\tau}
&=
\frac{\partial K}{\partial p_\mu},
\\
\frac{dp_\mu}{d\tau}
&=
-
\frac{\partial K}{\partial x^\mu}
+
F^\mu_{\rm diss}
+
\xi^\mu(\tau),
\end{align}

where

\[
F^\mu_{\rm diss}
\]

represents the dissipative contribution and

\[
\xi^\mu
\]

is a covariant stochastic force.

The precise statistical properties of

\[
\xi^\mu
\]

remain to be determined.

%%%%%%%%%%%%%%%%%%%%%%%%%%%%%%%%%%%%%%%%%%%%%%%%%%%%%%%%%%%%%%%%%%%%%%%%%%%%%%%%
\section{Memory Effects}
%%%%%%%%%%%%%%%%%%%%%%%%%%%%%%%%%%%%%%%%%%%%%%%%%%%%%%%%%%%%%%%%%%%%%%%%%%%%%%%%

If stochasticity originates from eliminating hidden degrees of freedom,
one expects

\begin{equation}
F^\mu_{\rm diss}
=
\int_{-\infty}^{\tau}
M^{\mu}{}_{\nu}
(\tau-\tau')
u^\nu(\tau')
\,d\tau',
\end{equation}

where

\[
u^\mu
=
\frac{dx^\mu}{d\tau}
\]

is the four-velocity.

The tensor

\[
M^{\mu}{}_{\nu}
\]

is the covariant memory kernel.

%%%%%%%%%%%%%%%%%%%%%%%%%%%%%%%%%%%%%%%%%%%%%%%%%%%%%%%%%%%%%%%%%%%%%%%%%%%%%%%%
\section{Covariant Noise}
%%%%%%%%%%%%%%%%%%%%%%%%%%%%%%%%%%%%%%%%%%%%%%%%%%%%%%%%%%%%%%%%%%%%%%%%%%%%%%%%

The stochastic force should satisfy

\begin{align}
\langle
\xi^\mu(\tau)
\rangle
&=
0,
\\
\langle
\xi^\mu(\tau)
\xi^\nu(\tau')
\rangle
&=
C^{\mu\nu}
(\tau-\tau').
\end{align}

The correlation tensor

\[
C^{\mu\nu}
\]

must transform covariantly.

Its detailed form depends upon the physical origin of the eliminated
degrees of freedom.

%%%%%%%%%%%%%%%%%%%%%%%%%%%%%%%%%%%%%%%%%%%%%%%%%%%%%%%%%%%%%%%%%%%%%%%%%%%%%%%%
\section{Fluctuation--Dissipation Relation}
%%%%%%%%%%%%%%%%%%%%%%%%%%%%%%%%%%%%%%%%%%%%%%%%%%%%%%%%%%%%%%%%%%%%%%%%%%%%%%%%

If both dissipation and fluctuations arise from the same microscopic
environment,

one expects a covariant fluctuation--dissipation theorem,

\begin{equation}
C^{\mu\nu}
(\tau)
=
\mathcal{F}
\!\left[
M^{\mu\nu}
(\tau)
\right],
\end{equation}

where

\[
\mathcal{F}
\]

denotes a model-dependent functional relation.

Deriving this relation from first principles is one of the principal
open problems of covariant stochastic dynamics.

%%%%%%%%%%%%%%%%%%%%%%%%%%%%%%%%%%%%%%%%%%%%%%%%%%%%%%%%%%%%%%%%%%%%%%%%%%%%%%%%
\section*{Critical Assessment}
%%%%%%%%%%%%%%%%%%%%%%%%%%%%%%%%%%%%%%%%%%%%%%%%%%%%%%%%%%%%%%%%%%%%%%%%%%%%%%%%

The equations proposed in this chapter should not be viewed as a
complete physical theory but rather as a set of structural principles.
They identify the mathematical ingredients that any covariant
stochastic extension of Hamiltonian mechanics must contain:
Hamiltonian evolution, invariant parametrization, covariance, memory,
fluctuations, and probability conservation.

The remaining challenge is to derive the explicit form of the
dissipation kernel and the stochastic correlation tensor from concrete
microscopic models. Depending on the underlying physics, these may arise
from electromagnetic fields, gravitational degrees of freedom, quantum
vacuum fluctuations, or many-body interactions. Consequently, the
framework developed here is intentionally general and can accommodate a
wide class of relativistic stochastic theories.




